\chapter{Conclusão}
\label{chapter:conclusao}

Esta dissertação apresentou o \textit{G-FShield}, uma arquitetura distribuída, orientada a eventos e observável para GRASP-FS em \acp{ids} aplicados a \acp{cps}. A implementação separa rankings e construção inicial, controle de vizinhanças, operadores de busca local, verificação e monitoramento; Apache Kafka integra os estágios, enquanto o backend e a interface permitem consultar sinais persistidos da execução.

A campanha formal respondeu à RQ1 com unidades ponta a ponta comuns. Foram avaliados 24 braços G-FShield, dois monólitos e um baseline com todas as características, cada qual em oito sementes independentes. RF--VND--IWSSR foi escolhido pela validação e obteve F1 macro mediano de 0,9448 no teste com 11 dos 51 preditores, redução de 78,4\%. O baseline obteve 0,9378, o Monólito~1 0,9397 e o Monólito~2 0,8411. A diferença diante do baseline é descritiva; nenhuma das 72 comparações exploratórias ajustadas por Holm apresentou $p<0{,}05$.

Para as RQ2 e RQ3, uma campanha causal adicional removeu o confundimento entre algoritmo e arquitetura: 30 sementes parearam implementações distribuída e monolítica da mesma sequência RF--VND--IWSSR, sob prazo de 2.700~s e teto agregado de seis CPUs e 12~GiB. O G-FShield avaliou 0,1074 candidato/s contra 0,0739 do monólito e sobrepôs praticamente 100\% da busca local à construção, contra 0\% no fluxo sequencial. Esses dois resultados ocorreram nos 30 pares. A maior vazão, contudo, não produziu menor latência de qualidade: o tempo mediano até F1 macro de validação 0,94 foi 940,1~s no distribuído e 701,3~s no monólito; a diferença pareada mediana foi $+422{,}7$~s, com $p=0{,}000795$ na direção contrária à hipótese. A AUC \textit{anytime} também favoreceu o monólito.

O paralelismo utilizou mediana de 4,076 núcleos, contra 1,004, e apresentou maior custo de CPU-h e GiB-h por candidato em todos os pares. Em compensação, a redução dimensional mediana foi 78,43\% no G-FShield e 76,47\% no monólito. Assim, a evidência demonstra a capacidade arquitetural de executar estágios concorrentemente e aumentar vazão bruta, mas não sustenta superioridade global: para uma única requisição, o monólito atingiu o limiar mais cedo e foi mais eficiente por recurso.

Quanto à RQ4, manifesto, comandos resolvidos, hashes, identificadores de campanha e execução, sementes, estados, motivos de parada, timestamps, mensagens, CSVs internos, logs, métricas finais e telemetria permitem reconstruir o que foi capturado. A aplicação também persiste eventos por semente e estágio. Não foram medidos completude ou custo da telemetria, perda e duplicação sob falhas, tempo de recuperação ou reprocessamento; resiliência, tolerância a falhas e escalabilidade permanecem objetivos de projeto.

O confundimento entre arquitetura e algoritmo permanece na comparação ampla de 27 braços, mas foi removido na campanha causal pareada. Suas principais ameaças remanescentes são as tentativas distribuídas repetidas após falhas de inicialização e o escopo restrito a um servidor, um corpus, um particionamento, J48, uma configuração de concorrência e 30 sementes. Os resultados não demonstram escalabilidade horizontal, tolerância a falhas ou generalização para outras cargas.

Como trabalhos futuros:

\begin{itemize}
  \item concluir a campanha confirmatória independente, já pré-registrada para as sementes 72--101, sobre o número de subconjuntos distintos com F1 macro de validação $\geq 0{,}945$;
  \item repetir a comparação causal com mais sementes e cargas simultâneas independentes;
  \item executar estudos de \textit{scale-out} que variem réplicas e concorrência e relatem vazão por custo;
  \item medir completude, \textit{overhead}, recuperação e reprocessamento da telemetria;
  \item repetir o protocolo com outros corpora, classificadores, particionamentos e hardware.
\end{itemize}

Em síntese, o G-FShield demonstra uma integração distribuída e rastreável de GRASP-FS e produz um conjunto experimental auditável. O experimento causal confirma que a decomposição permite sobreposição entre construção e busca local e maior vazão de candidatos, mas revela o preço dessa capacidade: maior consumo por candidato e pior tempo até F1 0,94 nesta carga. A contribuição arquitetural é, portanto, um compromisso mensurado, não uma alegação de superioridade universal ou prontidão operacional.
